\documentclass[11pt]{article}
\usepackage[margin=1in]{geometry}
\usepackage[T1]{fontenc}
\usepackage{lmodern}
\usepackage{microtype}
\usepackage{enumitem}
\usepackage{amsmath,amssymb}

\begin{document}

\section{System Architecture and Methods}
\label{sec:architecture-methods}

The implementation follows a three-stage prioritized-planning pipeline and
extends it with a hybrid central-plus-floor-server architecture. Algorithm~1
assigns jobs and constructs each AGV's task route; Algorithm~2 schedules those
routes in priority order; and Algorithm~3 shifts conflicted route windows to
the right. In the hybrid version, the central server handles assignment and
lift reservations, while floor servers plan route legs using local traffic
reservations. Function and source-line references below identify the
implementation in the repository.

\subsection{Algorithm 1: prioritized job assignment}

The assignment stage processes jobs in increasing desired pickup time
$t_j$, breaking ties by job identifier. For each candidate AGV $i$, the
fitness routine \texttt{fitness()} computes the predicted completion time
$ct_{ji}$, pickup earliness $et_{ji}$, delivery tardiness $dt_{ji}$, and
\[
  F_{ji}=ct_{ji}+\lambda(et_{ji}+dt_{ji}).
\]
The AGV with the smallest $F_{ji}$ is selected; ties are resolved by fixed AGV
priority. The assigned AGV's available time and current location are then
updated to its predicted completion time and the job's delivery node. For
each assigned job, the route is extended along shortest paths from the AGV's
current location to the pickup node and then to the delivery node, with
$\eta$ loading time at pickup and $\eta'$ unloading time at delivery. After
all jobs are assigned, each AGV's route is extended along a shortest path
back to the parking node $\pi$.

\paragraph{Pseudocode.}
\begin{verbatim}
for job j in sort(J, key=(t_j, job_id)):
    for AGV i:
        compute F_ji, ct_ji, et_ji, dt_ji from current (s_i, vt_i)
    choose i* = argmin_i (F_ji, priority_i)
    append shortest path s_i* -> p_j; add pickup dwell eta
    append shortest path p_j -> p'_j; add delivery dwell eta'
    update s_i* = p'_j and vt_i* = ct_ji*
for AGV i:
    append shortest path s_i -> pi
\end{verbatim}

The implementation is in \texttt{fitness()} (\texttt{agv\_core.py},
lines 149--157), the route-appending helper \texttt{\_append\_path()}
(lines 159--166), and \texttt{algorithm\_1()} (lines 168--196).

\subsection{Algorithm 2: prioritized conflict-free routing}

Given the routes from Algorithm~1, Algorithm~2 schedules AGVs in increasing
priority rank, i.e., $V_1$ through $V_{|V|}$ (in the reference experiment,
$V_1$ through $V_{20}$). For each route, \texttt{plan\_route()} constructs
node occupancy windows $[lt_m,rt_m]$ and edge traversal windows
$[rt_{m-1},lt_m]$. When a new node or edge window conflicts with a frozen
higher-priority reservation, it invokes Algorithm~3 to delay the affected
route prefix. After a route is planned, \texttt{freeze()} adds its ordinary
node windows to $W_n$ and its directed edge windows to $W'_e$. Track checks
consider reservations for both directions of the physical track. The parking
node is exempt from node conflicts.

\paragraph{Pseudocode.}
\begin{verbatim}
initialize node reservations W_n and directed edge reservations W'_e
for AGV i in ascending priority order V1, V2, ...:
    build node/edge windows along route R_i
    whenever a new window conflicts with W_n or W'_e:
        call Algorithm 3 to shift the affected prefix
    freeze ordinary node windows in W_n
    freeze edge windows in W'_e
\end{verbatim}

This stage is implemented by \texttt{algorithm\_2()} (\texttt{agv\_core.py},
lines 306--324), which calls \texttt{plan\_route()} (lines 269--295) and
\texttt{freeze()} (lines 297--304).

\subsection{Algorithm 3: sliding-window conflict elimination}

For a candidate occupancy interval $[lt,rt]$ and a set $W$ of fixed higher-
priority occupancy intervals, the function $f(lt,rt,W)$ returns $st$ and
$mt$. Here, $st$ is the smallest rightward shift that clears the candidate
window from all windows in $W$, padded by safety separation $\delta$; $mt$
is the available shift before the next fixed window would be reached. If the
candidate is already conflict-free, $st=0$.

When a conflict occurs at route visit $p$, \texttt{algorithm\_3()} sweeps
backward from that visit to the nearest preceding exempt node $q$ (the
parking node $\pi$ in the centralized case). It accumulates local shifts
along the node and edge windows, checks the running $mt^*$ bound, and repeats
the sweep until the segment is conflict-free or $k_{\max}$ passes have been
used. The resulting cumulative delay $st^*$ shifts the windows from $q+1$
through $p$ and extends the window at $q$ by the same amount. Thus, in the
centralized planner, the AGV waits at $\pi$ rather than stopping on an edge
or at an ordinary node. If the loop reaches $k_{\max}$ without convergence,
the fallback delays departure at the exempt node until the higher-priority
traffic is free.

\paragraph{Pseudocode.}
\begin{verbatim}
q = nearest exempt route node before conflict visit p
repeat for k = 1,...,k_max:
    st_star = 0; mt_star = infinity
    sweep m backward from p to q+1:
        (st, mt) = f(candidate node/edge window shifted by st_star, W)
        update mt_star and st_star
        stop this pass if st_star exceeds mt_star
    if the sweep completes without violating mt_star: break
if all k_max passes are exhausted:
    apply fallback delay at q until higher-priority traffic is free
shift windows q+1,...,p by st_star; extend the window at q by st_star
\end{verbatim}

The interval function $f()$ and the sliding procedure are implemented in
\texttt{agv\_core.py}, lines 198--216 and 218--260, respectively. Conflict
detection and calls to Algorithm~3 are in \texttt{plan\_route()} (lines
269--295).

\subsection{Hybrid central-plus-floor-server extension}

The hybrid planner first runs Algorithm~1 to determine job assignments and
complete routes. The helper \texttt{split\_legs()} then cuts each route at
every lift edge, producing maximal route segments that remain on one floor.
Each segment starts at the parking node or at a lift lobby and is submitted
to the corresponding \texttt{FloorServer}. A floor server maintains its own
local node-window set $W_n$ and edge-window set $W'_e$ and uses
\texttt{plan\_route()} (Algorithms~2 and~3) to plan each leg. When local
conflicts require a delay, the leg waits at its entry, which is the parking
node or the arrival lobby. Reservations are frozen locally after the leg is
planned.

The \texttt{CentralServer} owns reservations for lift tracks.
\texttt{book\_lift()} finds the earliest departure at or after the AGV is
ready for which the lift is available, respecting $\delta$ separation from
other trips on the same lift. Any lift-related delay is spent at the
departure lobby. The central process merges the planned floor legs and lift
windows into each complete route before independent verification.

Legs are submitted through a heap. With \texttt{order="fcfs"}, the heap key is
(release time, AGV priority, leg index), so a floor plans the earliest
released leg first, with priority used for ties. With
\texttt{order="priority"}, the key is (AGV priority, leg index, release
time), so every leg of $V_1$ is processed before legs of $V_2$, and so on.
This priority mode is used as a decomposition check: on the tested instances,
it reproduces the centralized-with-lobbies plans.

\paragraph{Pseudocode.}
\begin{verbatim}
run Algorithm 1 to assign jobs and construct each complete route
for AGV i:
    legs[i] = split_legs(route[i]) at every lift edge
push each first leg into heap using selected order key
while heap is not empty:
    pop (AGV i, leg k, release time)
    floor server plans leg with local W_n and W'_e
    if another floor leg follows:
        central server books the intervening lift trip
        submit the next leg at the arrival lobby after lift travel
merge all floor-leg and lift windows into complete AGV plans
\end{verbatim}

The hybrid functions are \texttt{split\_legs()} (\texttt{agv\_hybrid.py},
lines 47--56), \texttt{FloorServer.plan\_leg()} (lines 65--76),
\texttt{CentralServer.book\_lift()} (lines 89--99), and
\texttt{run\_hybrid()} (lines 106--165), which implements heap ordering,
leg dispatch, and merging. The floor graph, lift lengths, and optional
landing lobbies are defined by \texttt{build\_three\_floors()} in
\texttt{agv\_layouts.py}, lines 22--31 and 47--85.

\end{document}
